% v229: synchronized 49-item Japanese reviewer response; 77 PDF-verified anchors.
% Full Japanese translation of response_round1.tex (not a journal-submission file).
% New numbered Lemma 2.2 is cited in R1.11 (fresh-identity completion).
% Locator table is regenerated from the latest verified English marked-PDF response.
% Source freshness is enforced by scripts/sync_and_check.py.
\documentclass[10pt,a4paper]{ltjsarticle}
\usepackage{luatexja-fontspec}
% Use the TeX Live default Japanese fonts (portable across CI runners).
\usepackage{amsmath,amssymb,booktabs,longtable,array}
\usepackage{stmaryrd}
\newcommand{\NT}[1]{\llbracket #1\rrbracket}
\usepackage[margin=20mm]{geometry}
\usepackage[hidelinks,unicode]{hyperref}
\usepackage{bookmark}
\usepackage{xcolor}
\setlength{\parindent}{1em}
\setlength{\parskip}{3pt}
\setlength{\emergencystretch}{3em}
\newcommand{\Comment}[1]{\par\medskip\noindent\textbf{#1　査読コメント（引用部分の全訳）：}\ }
\newcommand{\Reply}{\par\noindent\textbf{著者回答：}\ }
\newcommand{\Location}[1]{\par\noindent\textbf{修正箇所：} #1\par}
\newcommand{\ResHead}[1]{\subsubsection*{#1}}
\newcommand{\Chcf}{\mathcal C_h^{\mathrm{cf}}}
\newcommand{\Chlin}{\mathcal C_h^{\mathrm{lin}}}
\newcommand{\Bh}{\mathsf{Rec}_h}
\newcommand{\Ah}{\mathsf{Learn}_h}
\newcommand{\RSh}{\mathsf{RS}_h}
\newcommand{\RS}{\mathsf{RS}}
\title{査読者への回答 --- TCS-D-26-00494\\\large 日本語全訳（著者確認用）}
\author{Takayuki Kuriyama}
\date{}
\begin{document}
\maketitle
\noindent\textbf{論文：}Distributional Learning of Context-Free Languages under Fixed Finite-Monoid Typing

\noindent\textbf{注記：}本書は英語版 \texttt{response/response\_round1.tex} の非提出用日本語訳である。引用された査読意見も省略部分の記号を保持して訳した。節・命題番号は改訂稿を指し、原稿の旧番号は「投稿時」と明記した。ページ・余白行番号表は英語版から自動的に同期しており、最新版の変更表示付きPDFで49項目・77参照の位置を再検証した。英語本文に変更があった場合は同期検査により未更新を検出する。

\input{japanese/parts/00_overview.tex}
\input{japanese/locator_table_JP.tex}
\input{japanese/parts/01_editor_r1_01_14.tex}
\input{japanese/parts/02_r1_15_30.tex}
\input{japanese/parts/03_r1_31_37.tex}
\input{japanese/parts/04_r2.tex}
\end{document}
